\subsection{Stafford型楕円包絡の中山作用と一般点の排除}
\label{subsec:stafford-exclusions}
標準Sklyanin以外にも、同じ楕円THCRへの全射を持つAS4包絡が存在する。
ここではDaviesの一次資料\cite[(6.1.3)--(6.1.4), Theorem 6.1.10]{DaviesThesis}
にあるStafford族を扱う。点加群の無限位数に関する結果は用いない。
パラメータ領域を
\[
 T=\{(\alpha,\beta,\gamma):\alpha+\beta+\gamma+\alpha\beta\gamma=0,
 \quad \alpha,\beta,\gamma\notin\{0,1,-1\}\}
\]
とし、$(i,j,k)=(1,2,3),(2,3,1),(3,1,2)$ に対して
\[
 f_{2i-1}=[x_0,x_i]-\alpha_i\{x_j,x_k\},\qquad
 f_{2i}=\{x_0,x_i\}-[x_j,x_k],\quad
 (\alpha_1,\alpha_2,\alpha_3)=(\alpha,\beta,\gamma)
\]
と置く。さらに
\[
 g_1=-x_0^2+x_1^2+x_2^2+x_3^2,\quad
 g_2=x_1^2+A x_2^2+C x_3^2,\quad
 A=\frac{1+\alpha}{1-\beta},\quad C=\frac{1-\alpha}{1+\gamma}.
\]
$S_\infty$ の関係空間は $\langle f_1,f_2,f_3,f_4,g_1,g_2\rangle$、
$S_{d,i}$ は $f_i$ を $d_1g_1+d_2g_2$ へ置換したものとする。
$i=1$ では $d\ne(1:0),(1:-1-\beta\gamma)$、
$i=2$ では $d\ne(1:\beta-1),(1:-1-\gamma)$ を仮定する。
この領域では元の族も標準Klein cocycle twistもnoetherian Koszul AS4整域である。
cocycle規約は前節と同じ $\deg(x_0,x_1,x_2,x_3)=(0,e_1,e_2,e_1+e_2)$、
$c((p,q),(r,s))=(-1)^{ps}$ であり、上添字 $c$ でtwistを表す。

\begin{theorem}[Stafford族の中山作用]
\label{thm:stafford-nakayama}
上の全許容パラメータに対して、順序 $(x_0,x_1,x_2,x_3)$ で
\[
\begin{array}{c|ccc}
 B&S_\infty&S_{d,1}&S_{d,2}\\\hline
 \mu_B&-I&\operatorname{diag}(-1,-1,1,1)&\operatorname{diag}(1,1,-1,-1)\\
 \mu_{B^c}&-I&-I&I
\end{array}
\]
である。従って $S_\infty,S_\infty^c,S_{d,1}^c$ は全体がH3包絡から除外される。
残る $S_{d,1},S_{d,2},S_{d,2}^c$ も各パラメータ領域の非空開稠密集合で
$W_1=0$ となる。循環置換で得られる他のStafford族にも対応する結論が成り立つ。
残る閉部分集合の全排除は主張しない。
\end{theorem}
\begin{proof}
各基底は既約である。実際 $1+\alpha\beta\ne0$ なので
$\gamma=-(\alpha+\beta)/(1+\alpha\beta)$ により $T$ は平面の開集合であり、
$d$ の二つの禁止切断を除く $\mathbb P^1$ 族も既約である。
有限群次数の四つの異なる指標で分解される $B_1$ 上で、
中山自己同型は符号群の全自己同型と可換するため対角である。
ここではgraded中山作用の内的曖昧さは次数零の単元、すなわちscalarのみである。

$S_{d,i}$ の群次数零の関係はただ一本 $Q_d=d_1g_1+d_2g_2$ である。
その四つの平方係数が全て非零となる稠密開集合では、
$\mu(x_j)=m_jx_j$ が $kQ_d$ を保つことから全 $m_j^2$ が等しい。
$S_\infty$ では平方関係のpencilを保つ。四つの比 $0,1,A,C$ は全て異なる。
例えば $A-C=\alpha$、$A-1=-\gamma C$、$C-1=\beta A$ である。
従ってその $2\times4$ 係数行列の全六Pl\"ucker座標は非零であり、
対角作用 $\operatorname{diag}(m_j^2)$ が行空間を保つなら
全 $m_i^2m_j^2$ は等しく、再び全 $m_j^2$ が等しい。
cocycle twistは平方係数の各列を非零符号倍するだけなので同じ議論が通る。

従って稠密開集合で $\mu=\lambda\epsilon$、$\epsilon$ は対角符号行列となる。
Koszul AS正則環の恒等式 $\operatorname{hdet}\mu=1$
\cite[Theorem 6.3]{RRZ} と、AS index $4$ による
$\operatorname{hdet}(\lambda I)=\lambda^4$ を使う。
$\epsilon^2=I$ より $\operatorname{hdet}\epsilon\in\{1,-1\}$ なので
$\lambda^8=1$。従って中山行列は有限集合に値を取る。

この有限性を一つの標本からの無条件な外挿と混同しないため、
族における正則性も確認する。二次関係行列から作るKoszul双対の
次数 $3,4$ 商は、全許容fiberで次元 $4,1$ が一定なので局所自由である。
AS4入力により双対Frobenius pairingは全fiberで非退化である。
最高次数の直線束を局所自明化すると、中山行列は二つのpairing行列の
積と逆行列で与えられ、従って正則であり自明化に依存しない。
既約な被約基底上のこの正則写像が稠密開集合で有限集合に値を取るので、
基底全体で一定である。

残る定数の決定は $(\alpha,\beta,\gamma)=(2,3,-5/7)$、$d=(1:4)$ で行う。
ここは全ての禁止値を避ける。検証器
\texttt{verify\_cubic\_stafford.py} は六つの関係空間からそれぞれ $R^\perp$ を計算し、
双対次数 $3,4$ 商と非退化pairingを有理数上で直接求める。
$\varphi$ を最高次数汎関数、$w_j$ を双対次数三の基底とし
$L_{ij}=\varphi(x_iw_j)$、$P_{ji}=\varphi(w_jx_i)$ と置くと、
$\mu_E=(L(P^{\mathsf T})^{-1})^{\mathsf T}$、
$\mu_B=-\mu_E^{-\mathsf T}$ が表の値を与える。
証明書にはpairing行列式と最高次数汎関数を保存してある。

$\mu_B=-I$ なら次数三では $\mu_B(G)=-G$ であり、
H3の必要条件 $\mu_B(G)=G$ と両立しない。
他の三族では同じ標本で $B_3,B_4$ の次元 $20,35$ と
$140\times20$ 行列の非零最大minorを直接計算した。
関係商が局所自由で中山行列が一定であるため、このminor非零条件は開であり、
各既約基底の開稠密部分で $W_1=0$ を得る。
循環置換は\cite[Lemma 6.1.11]{DaviesThesis}に従う。
\end{proof}

\paragraph{未解決の特殊部分への還元.}
残余三族の全H3候補は、それぞれ上の線型写像の全 $20$ 次minorが消える
閉部分集合に支持される。この条件はadmissibilityの十分条件ではない。
全パラメータのminor idealの簡約、特殊包絡の正規三次の分解、
商の内在的H3判定は依然として必要である。

\begin{lemma}[楕円商で消える三次は整域切断にならない]
\label{lem:elliptic-normal-sequence-cubic}
$A$ を次数一生成のgraded整域とし、二つの二次式による正則正規列
$(q_1,q_2)$ の商を $C=B(E,L,\sigma)$ とする。
ここで $E$ は滑らかな楕円曲線、$L$ は次数四のline bundle、
$\sigma=t_\tau$、$4\tau\ne0$ とする。
$G\in A_3$ が正規で、そのnormalizer $\theta$ が $q_1$ をscalar倍し、
$G$ の $C$ への像が零なら、$A/(G)$ は整域ではない。
\end{lemma}
\begin{proof}
まず $C$ に非零の正規一次元は存在しない。
そのような $\ell$ と可逆線型写像 $\psi:C_1\to C_1$ の関係
$\ell v=\psi(v)\ell$ を $H^0(E,L\otimes\sigma^*L)$ で評価する。
$\ell$ と $\sigma^*\ell$ の零点を除く稠密開集合上で、
全 $v\in H^0(E,L)$ の比から、$\psi$ の射影作用が
$L$ による埋込み上の $\sigma$ と一致する。
従って $\sigma^*L\simeq L$、すなわち $4\tau=0$ となり矛盾する。
ここではTHCRの全自己同型分類を必要としない。

$\overline A=A/(q_1)$ と置く。$\theta$ はこの商へ降り、
$G$ の像は $\ell q_2$、$\ell\in\overline A_1$ と書ける。
$q_2 v=\nu(v)q_2$ とすると、$q_2$ の右正則性による消去から
\[
 \ell u=\theta\nu^{-1}(u)\ell\quad(u\in\overline A_1)
\]
を得る。この式をさらに $C$ へ写す。
$\ell\ne0$ なら次数から $C$ における像も非零である。
$\theta\nu^{-1}$ の次数一の可逆行列が $C$ の関係を保つことも、
上の式と $C$ の整域性で消去すれば従う。
従ってこれは $C$ の正規一次元を与え、前段に反する。
ゆえに $G\in(q_1)$、すなわち $G=\ell_1q_1$ である。
$G\ne0$ なら両因子は非零で次数 $1,2$ を持つため、
$A/(G)$ に零因子が生じる。
\end{proof}

\begin{corollary}[通常のStafford族は位数六へ還元される]
\label{cor:stafford-order-six-only}
上の $S_{d,1},S_{d,2}$ の対応する楕円THCRを
$B(E,L,t_\tau)$ とし、$E$ が滑らか、$4\tau\ne0$ とする。
これらがH3包絡となるなら $\tau$ は位数六である。
従って無限位数および位数六以外のproper楕円部分は全て排除される。
この系はcocycle twistの非整域THCR商には適用していない。
\end{corollary}
\begin{proof}
Davies\cite[Theorem 6.1.10の証明]{DaviesThesis}で確認されるように、
楕円THCRへの核は二次式の正則正規列で生成される。
この核の次数二部分は $f_i$ と $g_1,g_2$ の残余一本で張られる。
定理\ref{thm:stafford-nakayama}の二つの通常の中山行列は、この二次元空間を
点ごとに固定する。従って列の第一元 $q_1$ は
H3三次のnormalizer $\mu_A^{-1}$ で固定される。
H3商は整域なので、前補題より $G$ の楕円商への像は非零である。

その像の正規関係を $E$ 上で評価すると、前補題と同様に
$\mu_A^{-1}$ の射影作用は $t_{3\tau}$ と一致する。
どちらの通常の中山行列も射影的には
$\operatorname{diag}(1,1,-1,-1)$ である。
これは $E$ 上固定点を持たない非自明な対合であり、
ある $s\in E[2]\setminus\{0\}$ による平行移動である。
実際その射影固定平面二つと $g_1=g_2=0$ の交わりは、
$A,C,A-C\ne0$ により空である。
従って $3\tau=s$、$6\tau=0$、$3\tau\ne0$。
位数二は $4\tau\ne0$ で除外され、位数は六に限られる。
\end{proof}

\begin{theorem}[非零候補が現れるStafford曲線も全排除される]
\label{thm:stafford-resonant-pencil-exclusion}
$a\notin\{0,1,-1,3,-3\}$ とし
\[
 (\alpha,\beta,\gamma)
 =\left(a,\frac{1-a}{a+3},\frac{a+1}{a-3}\right)
\]
とする。対応する楕円商は滑らかで $4\tau\ne0$ と仮定する。
$S_{d,1}$ の全許容 $d\in\mathbb P^1$ において $\dim W_1=2$ であり、
全ての非零 $G\in W_1$ が一次式と二次式の積に分解する。
従ってこの曲線上の $S_{d,1}$ は全てH3包絡から除外される。
この曲線の外の位数六部分全体を排除したという主張ではない。
\end{theorem}
\begin{proof}
次の式は $d$ に依存しない。
\[
\begin{split}
 Q_2&=2\bigl((a-3)(x_0^2+x_3^2)+(a+1)(-x_1^2+x_2^2)\bigr),\\
 Q_3&=2\bigl(-(a+3)(x_0^2+x_2^2)+(a-1)(x_1^2-x_3^2)\bigr),\\
 F_2&=x_2Q_2,\qquad F_3=x_3Q_3,\\
 H&=-16a x_1x_0+8(a-1)x_2x_3-8(a+1)x_3x_2.
\end{split}
\]
検証器\texttt{verify\_cubic\_stafford\_resonance.py}は $\mathbb Q(a)$ 上で、
二次関係 $f_2,\ldots,f_6$ だけを課した商において
\[
 F_jx_i=\epsilon_i x_iF_j\quad
 (j=2,3),\qquad (\epsilon_0,\epsilon_1,\epsilon_2,\epsilon_3)=(-1,-1,1,1)
\]
と、$g_1,g_2$ の中心性、$f_1$ の同じ $\epsilon$ による正規性を直接確認する。
従って全許容Stafford族でもこれらの恒等式が成り立つ。
特殊パラメータへの延長は、全fiberが14641型であることによる各次数の局所自由性を使う。
有理関数体でのrow reductionの分母を、未確認の特殊値で形式的に反転してはいない。

四つの平方係数を比較すると、$Q_2,Q_3$ のいずれも
$k(d_1g_1+d_2g_2)$ に属さない。従って $F_2,F_3$ は各fiberで非零であり、
異なる群次数に属するので独立である。これで $\dim W_1\ge2$ となる。

次に $W_1$ から楕円THCRへの写像は単射である。
実際、核の第一二次元として $g_1,g_2$ の非零の残余中心元 $q$ を取れる。
$q$ と $f_1$ は正則正規列を成す。
第一元の正則性は整域性、第二元はTHCRのHilbert級数との比較による。
補題\ref{lem:elliptic-normal-sequence-cubic}の証明により、
楕円商で消える $G\in W_1$ は $\ell q$ である。
$q$ は中心なので、消去すると $\ell$ もnormalizer $\mu_A^{-1}$ を持つ。
その楕円商への像は非零の正規一次元となり、$4\tau\ne0$ に反する。

非零の像があるので、前系の評価の議論から $\tau$ は位数六である。
同一normalizerを持つ非零三次切断の比は $\sigma$ 不変な有理関数であり、
その零因子も $\sigma$ 不変である。
自由な六重被覆 $\pi:E\to E/\langle\tau\rangle$ で降下すると、
次数十二の零因子は次数二の因子となる。
従ってそのような切断の空間は次数二のline bundleの切断空間へ入り、次元は高々二である。
単射性と先の独立性から $W_1=kF_2\oplus kF_3$ を得る。

最後に次の普遍恒等式を使う。
\[
\begin{split}
 &(s x_2+t x_3)
 \bigl((a-1)s^2Q_2+stH-(a+1)t^2Q_3\bigr)\\
 &\quad=s\bigl((a-1)s^2+(a+3)t^2\bigr)F_2
       -t\bigl((a-3)s^2+(a+1)t^2\bigr)F_3.
\end{split}
\]
検証器は $s^2t,st^2$ の二係数を五つの二次関係から厳密検証する。
右辺の二係数は $\mathbb P^1_{s:t}$ 上に共通零点を持たない。
$st\ne0$ では二つの二次式の係数行列式が $8$、
両端では $a-1,a+1\ne0$ を使う。
従って次数三の全射 $\mathbb P^1\to\mathbb P(W_1)$ を定める。
全ての非零 $G$ が表示された積に分解し、商に零因子が生じる。
\end{proof}
